% Page 026 — page imprimée 22
% Contenu seul : le préambule et la pagination sont gérés par meyrueis.tex

\begin{center}
{\bfseries Histoire de Samuel François}

\vspace{0.45em}

{\large\bfseries Pasteur de Meyrueis, nommé Administrateur\\
du Département de la Lozère}

\vspace{0.45em}

{\small\bfseries\itshape
D’après les comptes-rendus de la Société Populaire de Meyrueis\\
Archives Gévaudanaises : 1913. JB Belon P. 293 à 316}
\end{center}

\vspace{0.7em}

A cette époque Meyrueis avait bénéficié du ministère pastoral de Samuel François.
Écœuré sans doute, par les excès de la Révolution il avait été nommé chef du parti Girondin de
Lozère. Il soutenait un certain fédéralisme au point qu’il fut jeté en prison à Mende. Il réussit
à s’évader et les Montagnards se jetèrent à sa poursuite.

C’est ainsi qu’arriva à Meyrueis, le 6 germinal (27 mars 1794), le Montagnard Chateauneuf-
Randon qui, prenant la parole « avance, dit-il, avec regret, qu’il est convaincu qu’il existe
dans la commune des citoyens instruits du lieu de la retraite de Samuel François ; il démontre
avec fermeté le crime de la récelation (sic) et propose de faire une adresse au peuple pour
l’éclairer de son devoir. »

A ces mots la Société par un mouvement spontané se lève et jure de révéler tout ce qu’elle
pourra découvrir de relatif à la retraite de « cet homme pervers qui a tenté, mais en vain,
d’avilir la représentation nationale et d’égarer le peuple sur ses véritables intérêts ».

Samuel François reste introuvable ; Chateauneuf-Randon décide alors de faire incarcérer sa
femme, mère de plusieurs enfants. Son sort émeut quelques anciens paroissiens de son mari à
Meyrueis, et lors d’une séance du Comité, le 24 vendémiaire an III (16 Octobre 1794) l’un
d’entre eux prend la parole et s’exprime en ces termes :

\begin{quote}
\itshape
« Citoyens, je viens intéresser votre justice, votre humanité en faveur de la citoyenne
François, dont le civisme, depuis 1789, ne s’est jamais démenti. Vous avez été les témoins de
sa conduite pendant la Révolution. Vous connaissez ses malheurs et la pureté de ses
intentions ; vous êtes instruits des motifs qui l’ont privée de sa liberté.

Vous savez qu’elle est en arrestation par ordre du représentant du peuple Châteauneuf-
Randon, depuis plus de six mois. Vous n’ignorez pas que l’évasion de son mari fut la seule
cause de cette mesure de sûreté générale. Vous savez encore que cette citoyenne a une
famille qui ne vit que de son travail et que son absence a réduite à la dernière misère. Ces
motifs de justice et d’humanité basent une réclamation qui tend à ce que la Société se
prononce en faveur d’une citoyenne intéressante par ses malheurs et son civisme, et qu’elle
sollicite sa mise en liberté auprès du comité de Sûreté Générale de la Convention. Plusieurs
membres parlent de la citoyenne François et tous manifestent le même sentiment à son égard.

L’assemblée arrête à l’unanimité que le comité de Sûreté Générale de la Convention sera
invité d’ordonner la mise en liberté de la citoyenne Samuel François de Meyrueis, détenue
dans la maison de réclusion de Mende par ordre du représentant du peuple Châteauneuf-
Randon. »
\end{quote}
